\subsection{Point injectivity after cyclotomic base change}

\begin{statement}{C.5}{Proved}{Descent of the unique point-section conjugator}
Let $F=\Q(\mu_\infty)$, let $K/F$ be finitely generated, and let
$X/K$ be a hyperbolic curve. If $x,y\in X(K)$ have geometrically
conjugate point sections on the cyclotomic kernel of $G_K$, then $x=y$.
The conjugator is unique. No local section theorem is used.
\end{statement}
\begin{proof}
Choose generators $z_1,\ldots,z_r$ with $K=F(z_1,\ldots,z_r)$.
Adjoin to $\Q(z_1,\ldots,z_r)$ the finitely many coefficients needed
to define $X$, its smooth compactification and boundary, and $x,y$.
Call the resulting subfield of $K$ $k$. It is finitely generated over
$\Q$, and the inclusions
\[
 K=F(z_1,\ldots,z_r)\subseteq Fk\subseteq K
\]
prove the exact equality $K=k(\mu_\infty)$. In particular $K/k$ is
algebraic and Galois. We may use one algebraic closure $\bar k$ for
both fields. Its fixed-field correspondence gives
\[
 G_K=\operatorname{Gal}(\bar k/k(\mu_\infty))
     =H_k:=\ker\bigl(\chi_k:G_k\to\widehat\Z^{\times}\bigr).
 \tag{C.5.1}
\]
Indeed, an element of $G_k$ is in this kernel exactly when it fixes
every root of unity. Since $\mu_\infty\subset K$, the cyclotomic
kernel of $G_K$ itself is all of $G_K$.

Write $X_0/k$ for the chosen model, $P=\pi_1(X_{0,\bar k})$, and
$E_0=\pi_1(X_0)$. Geometric base change and the covering
correspondence identify the arithmetic group over $K$ with the
pullback
\[
 \pi_1(X)=E_0\times_{G_k}H_k.
 \tag{C.5.2}
\]
For the infinite algebraic extension this follows by taking the
inverse limit of the same identifications over its finite subextensions.
The geometric group is literally the same $P$ after the common
algebraic closure has been fixed; compare
\cite[Expos\'e IX, \S6]{AgentSGA1}.

Choose point-section representatives $s_x,s_y:G_k\to E_0$ and a
geometric conjugator $\gamma\in P$ satisfying
\[
 s_y(h)=\gamma s_x(h)\gamma^{-1}\qquad(h\in H_k).
 \tag{C.5.3}
\]
We use R15 precisely over the finitely generated field $k$, as
proved in LM6-F.7 \cite[F.7]{AgentLM6}. The group $s_x(H_k)$ fixes
the chosen lifted interior point above $x$, and R15 says this is its
only completed geometric closed pro-point. An element of
$C_P(s_x(H_k))$ preserves that singleton. The geometric deck action
on the fibre over an interior point is free, so
\[
 C_P(s_x(H_k))=1.
 \tag{C.5.4}
\]
Consequently the conjugator in (C.5.3) is unique. This application
does not extend the hypotheses of R15 from $k$ to $K$.

For $\sigma\in G_k$ put
$\gamma_\sigma=s_y(\sigma)\gamma s_x(\sigma)^{-1}$.
Its image in $G_k$ is $1$, so $\gamma_\sigma\in P$. Normality of
$H_k$ and (C.5.3) give, for $h\in H_k$,
\[
 \begin{split}
 \gamma_\sigma s_x(h)\gamma_\sigma^{-1}
 &=s_y(\sigma)\gamma s_x(\sigma^{-1}h\sigma)
                         \gamma^{-1}s_y(\sigma)^{-1}\\
 &=s_y(\sigma)s_y(\sigma^{-1}h\sigma)s_y(\sigma)^{-1}
 =s_y(h).
 \end{split}
\]
Uniqueness implies $\gamma_\sigma=\gamma$. Thus (C.5.3) holds
for every $h\in G_k$.

We include the last point-injectivity step. For a finitely generated
field $l/\Q$, the map
$l^\times\to\varprojlim_n l^\times/(l^\times)^n$ is injective.
To prove this, let $b$ belong to every $(l^\times)^n$, and let
$l_0$ be the algebraic closure of $\Q$ in $l$, a number field.
On a normal proper model of $l/l_0$, every prime-divisor valuation
of $b$ is divisible by every positive integer and hence zero.
Normality makes $b$ and $b^{-1}$ regular everywhere; properness
makes them constants. Thus $b\in l_0^\times$. Every root of $b$
lying in $l$ is algebraic over $l_0$ and also lies in $l_0$.
At the finite primes of $l_0$ its valuations are zero, so all these
roots are units. The unit theorem and
$\bigcap_{n\ge1}nM=0$ for a finitely generated abelian group $M$
give $b=1$. If $l=l_0$, simply start with these finite-prime
valuations. The group assertion follows by separating its free
part and then choosing $n$ divisible by its finite torsion exponent.

In particular, for the tripod the morphism $t:U_k\to\mathbb G_{m,k}$
sends the conjugacy of sections to
$\kappa_k(x)=\kappa_k(y)$; injectivity just proved gives $x=y$.
For a general genus-zero hyperbolic curve, make a finite extension
of $k$ splitting three boundary points. Its compactification has a
rational point $x$ and is a projective line; send the three boundary
points to $0,1,\infty$. The resulting coordinate is an injective
map on points into $\mathbb G_m$, and the same argument applies.

If the compactification $\bar X_0$ has positive genus, the map
$j_x:\bar X_0\to J=\operatorname{Pic}^0(\bar X_0)$,
$z\mapsto[z-x]$, sends the two point sections to the Kummer
classes of $0$ and $[y-x]$. Their equality implies
$[y-x]\in\bigcap_n nJ(k)$. The group $J(k)$ is finitely generated:
apply Lang--N\'eron over the number field of constants and then
Mordell--Weil to its trace; the precise Lang--N\'eron statement is
\cite[Theorem 2.1]{AgentConradTrace}. Thus $[y-x]=0$.
A principal divisor $y-x$ with $x\ne y$ would give a degree-one
map $\bar X_0\to\Pj^1$, contradicting positive genus. Hence
$x=y$ in this case as well.
\end{proof}

\subsection{The first separating lower-central stage}

\begin{statement}{C.6}{Proved}{A degree-$27$ Heisenberg cover}
Put $\xi=\zeta_6$ and $\eta=\xi^{-1}$ in $F$.
The point sections at $\xi,\eta\in U(F)$ coincide after geometric
abelianization and differ after quotienting by
$\Gamma_3\Pi=\overline{[\Pi,\overline{[\Pi,\Pi]}]}$.
An explicit geometrically connected finite \'etale Galois cover of
degree $27$, with constant group $\operatorname{UT}_3(\mathbf F_3)$,
distinguishes them. Thus class two is the first separating stage
of the lower-central filtration. No local theorem is used.
\end{statement}
\begin{proof}
Let $a=\zeta_{18}$, $\omega=a^6=\zeta_3$, and $\lambda=1-\omega$.
The Fermat cubic and its chosen origin are
\[
 E:w^3+v^3=1,\qquad O=(1:-1:0).
\]
Its Weierstrass model and the inverse formulas are
\[
 E:Y^2=X^3-432,\qquad
 X=\frac{12}{w+v},\quad Y=\frac{36(w-v)}{w+v},
 \qquad
 w=\frac{36+Y}{6X},\quad v=\frac{36-Y}{6X}.
 \tag{C.6.1}
\]
The displayed rational maps extend to inverse isomorphisms of
smooth proper curves. Write $[\omega](X,Y)=(\omega X,Y)$.
This is a group automorphism fixing $O$, and
$[\lambda]=[1]-[\omega]$ is its difference from the identity.

\emph{The arithmetic annihilator.}
We claim
\[
 E(F)[3^\infty]\subseteq E[\lambda^3].
 \tag{C.6.2}
\]
The ring $\mathcal O=\Z_3[\omega]$ is the discrete valuation ring
of $\Q_3(\omega)$: $\omega-1$ satisfies the Eisenstein polynomial
$T^2+3T+3$. The Tate module $T_3E$ is a torsion-free finite
$\mathcal O$-module, free of rank one since it has rank two over
$\Z_3$. Faithfulness of the endomorphism action, the torsion rank,
and the good-reduction Frobenius comparison are the standard
statements in \cite[II, Propositions 6.3, 6.4, 6.7; V,
Propositions 8.1, 8.3]{AgentMilneEC}.

At $7$ the curve has good reduction and $\#E(\mathbf F_7)=9$.
Indeed $X^3-432=X^3+2$ modulo $7$; the values for $X=0$, for
the three nonzero cubes equal to $1$, and for the three equal to
$-1$ contribute respectively $2,0,6$ affine points, followed by
$O$. Frobenius therefore has polynomial $T^2+T+7$.
It belongs to $G_{\Q(\omega)}$ because $7$ splits in
$\Q(\omega)$, so it commutes with $\mathcal O$ and acts by
$\alpha=1+3\omega$ or $\bar\alpha$. Complex conjugation acts
semilinearly and exchanges these scalars. Their commutator acts
by $\bar\alpha/\alpha$ or its inverse. Every commutator fixes
$F$, since all cyclotomic extensions of $\Q$ are abelian.
As $\alpha$ is a $3$-adic unit, an $F$-rational $3$-primary
torsion point is killed by
\[
 \alpha-\bar\alpha=3+6\omega=-(1-\omega)^3=-\lambda^3.
\]
This proves (C.6.2).

\emph{The required torsion points.}
Put $P=(a,a^{-1})$ in Fermat coordinates; it lies on $E$ because
$a^3+a^{-3}=1$. Direct application of the chord-and-tangent law
in (C.6.1) gives
\[
 [3]P=(0,24a^3-12)\ne O,\qquad
 [\lambda]^2P=(0,12-24a^3),\qquad [\lambda]^3P=O.
 \tag{C.6.3}
\]
Here is an explicit algebraic prescription for the first equality.
Set $X_1=12/(a+a^{-1})$, $Y_1=36(a-a^{-1})/(a+a^{-1})$, and
\[
 \begin{gathered}
 m_2=\frac{3X_1^2}{2Y_1},\quad
 X_2=m_2^2-2X_1,\quad Y_2=m_2(X_1-X_2)-Y_1,\\
 m_3=\frac{Y_2-Y_1}{X_2-X_1},\quad
 X_3=m_3^2-X_1-X_2,\quad Y_3=m_3(X_1-X_3)-Y_1.
 \end{gathered}
\]
Reduction by $a^6-a^3+1=0$ gives $X_3=0$ and
$Y_3=24a^3-12$. These denominators are nonzero: with
$d=a+a^{-1}$ we have $d^3=3d+1$, and
$X_2-X_1=-3X_1(X_1^3-1728)/(4Y_1^2)\ne0$.
The other equalities in (C.6.3) follow from
$\lambda^2=-3\omega$ and the fact that $[3]P$, whose
$X$-coordinate is zero, is fixed by $\omega$.

Let $\phi:E\to\Pj^1$ be $\phi(w,v)=w^3$. It has degree nine,
with deck group $\mu_3^2$ acting by separate multiplication of
$w,v$, and branches only at $0,1,\infty$. Its translation subgroup
is $E[\lambda]$. One can check this directly: the invariant
differential $dw/v^2$ is multiplied by $\omega^{i+j}$ under
$(w,v)\mapsto(\omega^i w,\omega^j v)$. The kernel consists of
three translations; they commute with $[\omega]$, so their
translation points are killed by $1-\omega$. Conversely
$\ker[\lambda]=\{O,(0,\pm12\sqrt{-3})\}$ has order three.
The diagonal scalings form the subgroup fixing $O$.
The swap $(w,v)\mapsto(v,w)$ is negation, as (C.6.1) shows.

The nine points above $\eta$ are
\[
 Q_{ij}=(\omega^i a^{-1},\omega^j a),\qquad 0\le i,j<3.
\]
Each has the form $T-[\omega]^rP$ with $T\in E[\lambda]$.
Since $[3]P$ is fixed by $\omega$ and has order three,
\[
 [3](Q_{ij}-P)=-([\omega]^r+1)[3]P=[3]P\ne O.
 \tag{C.6.4}
\]
It follows from $\lambda^2=-3\omega$ that
$[\lambda]^2(Q_{ij}-P)\ne O$. Equation (C.6.3) and
$[\lambda]T=0$ also give $[\lambda]^3(Q_{ij}-P)=O$.

\emph{The cover and its constant deck group.}
Define
\[
 M:E_F\longrightarrow E_F,\quad R\longmapsto[\lambda]R+P,
 \qquad f=\phi\circ M,
 \qquad W=f^{-1}(U_F).
 \tag{C.6.5}
\]
The isogeny $[\lambda]$ has degree three, by its kernel computed
above, and is \'etale in characteristic zero. Hence $W\to U_F$
is finite \'etale of degree $27$. It is geometrically connected,
being a nonempty open of the geometrically integral curve $E$.

For $b\in E[3]=E[\lambda^2]$, translation $t_b$ is a deck
transformation, since $\lambda b\in E[\lambda]$. So is
\[
 h(R)=[\omega]R-P:
 \quad M(hR)=[\omega]M(R),\quad \phi([\omega]Q)=\phi(Q).
\]
They satisfy
\[
 h^3=1,\qquad h t_bh^{-1}=t_{\omega b}.
\]
They give $27$ different automorphisms, and thus the full deck
group of the degree-$27$ map:
\[
 H_{27}=E[3]\rtimes\langle h\rangle.
\]
All are defined over $F$. Besides $P\in E(\Q(a))$, the complete
list of nonzero $3$-torsion points is
\[
 (0,\pm12\sqrt{-3}),\qquad (12\omega^i,\pm36),\quad i=0,1,2,
\]
as substitution or the third division polynomial confirms.
Thus the group scheme of deck transformations is constant over $F$.
The commutators are
$[h,t_b]=t_{(\omega-1)b}$, whose image is exactly $E[\lambda]$.
They are central, and the action of $\omega$ on $E[3]$ is
nontrivial. The group has class two. Every element has cube one,
using $1+\omega+\omega^2=0$. Choose $b$ with
$(\omega-1)b\ne0$. Then $h,t_b,[h,t_b]$ satisfy the usual
presentation of $\operatorname{UT}_3(\mathbf F_3)$; the map sending
them to $I+e_{12},I+e_{23},I+e_{13}$ is an isomorphism.

\emph{The two fibres.}
The point $R=O$ lies above $\xi$. An $F$-point $R$ above $\eta$
would satisfy $[\lambda]R=Q_{ij}-P$ for some $i,j$.
The right side is killed by $\lambda^3$, so $R$ is killed by
$\lambda^4=9\omega^2$ and is $3$-primary torsion. By (C.6.2),
$\lambda^3R=0$. Applying $\lambda^2$ to the preceding equation
contradicts (C.6.4). Hence the fibre above $\eta$ has no $F$-point.

A finite \'etale fibre is the finite $G_F$-set obtained by pulling
back its monodromy along the point section. Geometric conjugacy
of sections gives isomorphic such sets, preserving their fixed
points. Here the constant $H_{27}$-torsor gives an arithmetic
quotient $\pi_1(U_F)\to H_{27}$, surjective on the geometric
group. It kills $\Gamma_3\Pi$, since $H_{27}$ has class two.
The different fixed-point sets therefore distinguish the sections
already in that quotient.

Finally, $\Pi^{\mathrm{ab}}=\widehat\Z(1)^2$, with coordinates
induced by $t$ and $1-t$. The point classes are
$(\kappa_F(\xi),\kappa_F(1-\xi))$ and
$(\kappa_F(\eta),\kappa_F(1-\eta))$.
Here $1-\xi=\eta$, $1-\eta=\xi$, and every root of unity has
compatible roots of all orders in $F$. Both classes are zero.
This proves the first-stage assertion.
\end{proof}

The exact program \texttt{code/cyclotomic\_class2\_check.py}
recomputes (C.6.3)--(C.6.4), all nine differences, and the full group
law. The proof retains the annihilator $\lambda^3$ in the
Frobenius-$7$ argument of the earlier \texttt{CYCLOTOMIC\_KERNEL}
return, giving degree $27$ in place of its degree $81$.
The assertion concerns the indexed lower-central filtration; no least
element in the partially ordered set of all separating finite
quotients is asserted.

\subsection{Continuous semilinear actions and Conjecture B-prime}

Let $\mathcal U=\varprojlim_i U_i$ be a geometric universal
pro-cover of a hyperbolic curve over $\bar K$, and put
$\Omega=\bigcup_i\bar K(U_i)$. Its completed geometric closed
pro-points are compatible geometric closed points of the smooth
proper compactifications of the finite levels. This set excludes
the generic scheme point. A semilinear automorphism acts on
$\bar K$ by the specified automorphism of that field and on the
pro-curve by compatible algebraic maps.

Give the group of such automorphisms the topology of pointwise
convergence on the discrete set $\Omega$: the pointwise stabilizer
of each finite subset of $\Omega$ is an identity neighbourhood.
An action of $G_K$ is continuous for this topology precisely when
every $f\in\Omega$ has an open stabilizer. The formulation of
Conjecture B$'$ used here is that every such continuous semilinear
$G_K$-action, for $K/\Q$ finitely generated, has a fixed completed
geometric closed pro-point, with the corresponding
cyclotomic-kernel strengthening. Conjecture B imposes in addition
a normalized profinite deck group with hyperbolic curve quotient.
These are existence assertions whose continuous arithmetic content
is retained; we do not assert their general truth here. The reading
of scan 63 removes the normalization hypothesis and retains the
continuity indicated on scan 59 \cite[scans 59, 63]{Prompt2026}.

\begin{statement}{C.7}{Proved}{A continuous B-prime action acquires a finite quotient}
Every action just defined normalizes the deck group of a finite
hyperbolic covering level descended over $K$. It is the action of a
continuous section for that quotient. The level can be chosen to
cover any specified initial level. Consequently B$'$ and B over
finite extensions have the same fixed-point existence assertion,
including the cyclotomic-kernel version supplied by R15.
\end{statement}
\begin{proof}
Fix an initial field $L_0=\bar K(f_1,\ldots,f_m)$. Each orbit
$G_K f_i$ is finite, because its stabilizer is open in the compact
group $G_K$. Let $L$ be the field generated over $\bar K$ by all
these finitely many orbit elements. Then $L\supseteq L_0$ and
$L$ is $G_K$-stable. Its generators belong to a common finite
covering field in $\Omega$; hence $L/L_0$ is a finite extension
unramified over the initial open curve. It defines a connected
finite \'etale cover $Y$ of that curve. Its deck subgroup
$H=\operatorname{Aut}_L(\Omega)$ satisfies
\[
 gHg^{-1}=\operatorname{Aut}_{gL}(\Omega)=H.
 \tag{C.7.1}
\]

The action descends to the actual open curve $Y$. More explicitly,
factor each resulting map through a common finite Galois covering
level; invariance under $H$ makes it descend to $Y$. Equivalently,
the field map gives an isomorphism of its smooth proper
compactification; every geometric point of $Y$ lifts to
$\mathcal U$, so the pro-curve action sends it into $Y$. Applying
the same argument to the inverse proves that the boundary is
preserved.

Choose a finite coefficient extension defining this marked curve
and all the chosen generators. Intersect the generators' open
stabilizers with its Galois group and take the normal core. The
result is $G_{K'}$ for a finite Galois extension $K'/K$; it fixes
the generators and acts only on their scalar coefficients.
The remaining action of $\operatorname{Gal}(K'/K)$ is an effective
semilinear descent datum for the marked curve. For example,
the line bundle $\omega_{\bar Y}(D)^{\otimes3}$ is very ample
for a hyperbolic marked curve. Descend its finite-dimensional
section space and its homogeneous ideal by finite Galois descent.
This constructs a marked curve $(\bar Y_0,D_0)$ over $K$, with
$Y_0=\bar Y_0\setminus D_0$ and $(Y_0)_{\bar K}\simeq Y$.

The given action on $\Omega$ fixes the descended function field
$K(Y_0)$. The universal covering correspondence therefore places
it in
\[
 1\longrightarrow H\longrightarrow\pi_1(Y_0)
   \longrightarrow G_K\longrightarrow1,
\]
as a splitting. It is continuous for the Krull topology because
the stabilizer of every finite set of functions is open. Conversely,
a continuous arithmetic section has exactly this open-stabilizer
property on every finite covering field.

Passing to $Y$ does not alter the pro-cover or its completed
closed pro-points: the covering levels above $Y$ form a cofinal
system. Thus the fixed-point questions use the same action and
the same set. This proves the B$'$/B comparison. Once a quotient
has descended, R15 identifies a nonempty cyclotomic fixed set
with the same singleton fixed by the whole $G_K$. This is the
precise use of the supplied R16, LM6-F.9
\cite[F.9]{AgentLM6}.
\end{proof}

\subsection{An abstract action with empty cyclotomic fixed set}

\begin{statement}{C.8}{Proved}{The local cyclotomic Kummer kernel}
Let $L/\Q_p$ be finite and
$H=\ker(\chi:G_L\to\widehat\Z^{\times})$. If
$b\in L^\times$ satisfies
$\kappa_p(b)|_H=0$ in $H^1(H,\Z_p(1))$, then $b$ is a root
of unity. No local section theorem is used.
\end{statement}
\begin{proof}
Choose compatible $p$-power roots of $b$ and their continuous
Kummer cocycle $c:G_L\to\Z_p(1)$. Since $H$ acts trivially on
the coefficients, every coboundary on $H$ is zero. Thus $c|_H=0$
as a map, and $c$ descends to a continuous cocycle on the compact
abelian quotient $A=\chi(G_L)$.
Choose $a_0\in A$ with $\chi_p(a_0)\ne1$; a finite extension
of $\Q_p$ cannot contain all $p$-power roots of unity.
The cocycle equation and commutativity give
\[
 (\chi_p(a_0)-1)c(a)=(\chi_p(a)-1)c(a_0)\qquad(a\in A).
\]
The right side is a coboundary. Write the nonzero scalar on the
left as $p^m$ times a $\Z_p$-unit. It follows that
$p^m\kappa_p(b)=0$ after inflation to $G_L$, and hence
$b^{p^m}\in(L^\times)^{p^n}$ for every $n$.

If the residue field has $q$ elements, then
\[
 L^\times=\varpi^{\Z}\times\mu_{q-1}\times U^1,
 \qquad U^1=1+\mathfrak m_L,
 \qquad \bigcap_{n\ge1}(L^\times)^{p^n}=\mu_{q-1}.
\]
For the last equality, infinite $p$-divisibility first kills the
valuation. Remove the prime-to-$p$ root of unity, supplied by
Hensel's lemma. Every finite quotient
$U^1/(1+\mathfrak m_L^r)$ is a finite $p$-group; an element
divisible by every $p^n$ has trivial image in each. These
quotients separate points, proving the equality. Therefore
$b^{p^m}$, and consequently $b$, is a root of unity.
\end{proof}

\begin{statement}{C.9}{Proved}{An explicit ultralimit with a continuous local restriction}
There is an abstract, discontinuous section
$s:G_\Q\to E_U$ whose restriction to a fixed decomposition
subgroup $G_{\Q_3}$ is geometrically conjugate to the point
section at
\[
 a=1+\sum_{r=1}^{\infty}3^{r!}\in U(\Q_3).
 \tag{C.9.1}
\]
It induces a semilinear action on the original universal pro-cover,
normalizing its original geometric deck group. The stabilizer of
$\sqrt t$ is nonopen.
\end{statement}
\begin{proof}
First, $a$ is transcendental over $\Q$. Put
$A_N=1+\sum_{r=1}^N3^{r!}$. Then
\[
 v_3(a-A_N)=(N+1)!,\qquad 1<A_N<2\cdot3^{N!}.
\]
For a nonzero $P\in\Z[T]$ of degree $d$ with $P(a)=0$,
the nonzero integers $P(A_N)$, for all sufficiently large $N$,
would be divisible by $3^{(N+1)!}$, while their absolute values
are at most $C(2\cdot3^{N!})^d$ for a fixed $C$. This is
impossible. Also $v_3(a)=0$ and $v_3(a-1)=1$.

Recursively choose primes $q_n\ne3$, each dividing none of the
earlier integers $x_i$, and solve
\[
 x_n\equiv a\pmod{3^n},\qquad
 x_n\equiv q_n\pmod{q_n^2},\qquad
 x_n\equiv1\pmod{q_i}\quad(i<n).
 \tag{C.9.2}
\]
The least allowed prime and the least positive CRT representative
make this an explicit sequence. The representative is at least
two. The modulus $3^n$ only needs the terms with $r!<n$ in
(C.9.1). We obtain
\[
 x_n\longrightarrow a\quad\hbox{in }\Q_3,
 \qquad v_{q_i}(x_j)=\delta_{ij}.
\]
In particular the square classes $[x_n]$ are independent.

Choose continuous representatives $s_n:G_\Q\to E_U$ for the
point sections at $x_n$. For a nonprincipal ultrafilter
$\mathfrak u$ on $\mathbb N$ define
\[
 s(g)=\lim_{n\to\mathfrak u}s_n(g)\qquad(g\in G_\Q).
 \tag{C.9.3}
\]
Compact Hausdorffness of $E_U$ gives the limit; continuity of
multiplication gives $s(gh)=s(g)s(h)$, and the projection to
$G_\Q$ gives $\operatorname{pr}s(g)=g$.

We verify its precise local restriction. Fix a finite arithmetic
quotient $\varphi:E_U\twoheadrightarrow D$ and put
$P=\varphi(\Pi)$. The associated finite \'etale Galois cover
$W\to U$ has constant quotient
$W/P=U\times\operatorname{Spec}k_0$, with
$\operatorname{Gal}(k_0/\Q)=D/P$
\cite[Expos\'e IX, \S6]{AgentSGA1}.
For $b\in U(\Q_3)$ sufficiently close to $a$, its fibre is
equivariantly isomorphic to the fibre at $a$, both for the
Galois action and for the deck action of $D$, with the identity
on this constant quotient. To see the needed uniformity, take a
finite Galois extension of $\Q_3$ splitting $W_a$. The \'etale
inverse-function form of Hensel's lemma gives an analytic sheet
through each of its finitely many points on a sufficiently small
ball about $a$. Choose one ball for all sheets and all the
finitely many deck and field automorphisms. Uniqueness of the
sheets proves equivariance. Their constant coordinates retain
their original values.

With compatible paths, an equivariant torsor isomorphism inducing
the identity modulo $P$ is multiplication by an element of $P$.
It follows that, for all sufficiently large $n$,
$\varphi s_n|_{G_{\Q_3}}$ is $P$-conjugate to
$\varphi s_a$. This orbit is a finite closed subset of the
pointwise function space $D^{G_{\Q_3}}$, so its ultralimit
$\varphi s|_{G_{\Q_3}}$ belongs to the same orbit.

The elements $\gamma\in\Pi$ implementing this conjugacy at a
given finite quotient form a nonempty closed subset. Any finitely
many conditions are implied by one condition at a common finite
refinement. Compactness of $\Pi$ therefore gives one $\gamma$
such that
\[
 \gamma s|_{G_{\Q_3}}\gamma^{-1}=s_a.
 \tag{C.9.4}
\]
Here the local arithmetic group is identified with
$E_U\times_{G_\Q}G_{\Q_3}$ by characteristic-zero geometric
base change. In particular the local restriction of $s$ is continuous.

For discontinuity globally, let $\tau:E_U\to\mathbf F_2$
be the character of $z^2=t$. Then
$\tau s_n=\kappa_2(x_n)$. Independence in (C.9.2) gives a
continuous surjection
\[
 \kappa:G_\Q\twoheadrightarrow\prod_{n\ge1}\mathbf F_2,
 \qquad g\longmapsto(\kappa_2(x_n)(g))_n:
\]
every finite projection is onto by Kummer theory and the compact
image is closed. The ultrafilter functional
$\epsilon((e_n))=\lim_{n\to\mathfrak u}e_n$ is additive.
It kills the dense direct sum and takes the constant-one vector
to one; hence it is discontinuous. A continuous surjection of
compact Hausdorff spaces is a quotient map, so
$\epsilon\kappa$ is discontinuous. By (C.9.3),
\[
 \tau s=\epsilon\kappa.
\]
This proves discontinuity of $s$ and identifies the nonopen
stabilizer of the function $\sqrt t$ as $\ker(\epsilon\kappa)$.

Each element of $E_U$ acts semilinearly on the universal pro-cover
and extends to its proper normalization levels. Formula (C.9.3)
therefore gives an abstract semilinear action on their inverse
system. It normalizes $\Pi\triangleleft E_U$. This is the
construction in AP-D.12, now with all its finite-level geometric
conjugacy steps checked \cite[D.12]{AgentAP}.
\end{proof}

\begin{statement}{C.10}{Refuted}{The universal fixed-point assertion for abstract actions}
The abstract action of C.9 has no completed geometric closed
pro-point fixed by $G_{\Q(\mu_\infty)}$. In particular, the
unique-fixed-point assertion fails for unrestricted abstract
semilinear actions even when the original deck group is normalized.
This statement uses no local section theorem.
\end{statement}
\begin{proof}
Suppose $\widetilde x$ were such a fixed pro-point, and let
$x\in\Pj^1(\bar\Q)$ be its image at the base compactification.
Use the embedding $\bar\Q\hookrightarrow\overline{\Q}_3$
underlying (C.9.4), and choose a finite $L/\Q_3$ containing $x$.
Conjugate the action and $\widetilde x$ by the geometric
$\gamma$ in (C.9.4). The base point $x$ is unchanged and now
$s|_{G_{\Q_3}}=s_a$. Put
\[
 H=G_L\cap G_{\Q(\mu_\infty)}
   =\ker(\chi:G_L\to\widehat\Z^{\times}).
\]
Then $s_a(H)$ fixes the transported pro-point over $x$.

Use the following morphism $f:U\to\mathbb G_m$:
\[
 \begin{array}{c|c|c}
 x&f(t)&f(x)\\\hline
 x\in U(\bar\Q)&t&x\\
 0&1-t&1\\
 1&t&1\\
 \infty&t/(t-1)&1
 \end{array}
\]
In every case $f$ is regular and nonzero at $x$. Thus the covers
$z^{3^n}=f(t)$ extend \'etale across $x$; at a cusp their
monodromy kills its inertia because $\operatorname{ord}_x(f)=0$.
The compatible fibre coordinates of $\widetilde x$ send
$s_a(H)$ into the point decomposition group at $f(x)$.
The point section at $a$ gives the Kummer class of $f(a)$ in
these same covers. Consequently
\[
 \kappa_3\bigl(f(a)/f(x)\bigr)|_H=0
                 \quad\hbox{in }H^1(H,\Z_3(1)).
 \tag{C.10.1}
\]
This can also be read directly by evaluating the compatible
$3^n$-th roots at the two fibres. Changing roots changes the
cocycle by a coboundary, and every such coboundary on $H$ is zero.
This argument uses the continuous local restriction (C.9.4),
without assuming continuity of the original global action.

C.8 makes $f(a)/f(x)$ a root of unity. Both $f(x)$ and every
root of unity are algebraic over $\Q$, so $f(a)$ is algebraic.
Each displayed $f$ is an invertible fractional-linear function
over $\Q$; hence $a$ would be algebraic. This contradicts C.9
and proves the empty fixed set, checking AP-D.14
\cite[D.13--D.14]{AgentAP}.
\end{proof}

The examples in the earlier returns have the following exact scopes.
LM6-F.41 constructs a discontinuous action with one fixed point.
LM6-F.47 works over $\Q(i)$ and remains discontinuous on every
open subgroup of the cyclotomic kernel, already on
$H_1(U_{\bar K},\mathbf F_3)$; its fixed point is still the chosen
lift of $1/2$, and it admits no continuously descended connected
pro-\'etale finite-type quotient on those subgroups
\cite[F.41, F.47]{AgentLM6}. AP-D.11 gives an abstract action
normalizing no finite-quotient deck group and having one fixed
cusp \cite[D.11]{AgentAP}. C.9--C.10 supply the empty fixed set.
The comparison with C.7 is the inclusion of continuous actions
into abstract actions, with the explicit nonopen stabilizer
$\ker(\epsilon\kappa)$ witnessing strictness. These constructions
do not supply a counterexample to the continuous assertions A1.
The program \texttt{code/local\_ultralimit\_check.py} verifies the
first twelve CRT terms and all $144$ valuation identities. The
ultrafilter choice is an existence argument, not an algorithm
for evaluating the resulting abstract section.
